%% kapa: exam=1 round=32 year=2021 period=1 form=A subject=economics subject_ko=경제학원론 seq=6 number=46
%% source: Q-Net 감정평가사 기출문제 (cst00309) artlSeq=5208397 file=extracted/1008/5208397/A형/2021년도 제32회 감정평가사 1차 1교시 A형.pdf
%% conversion: kordoc 4.24.51 -> tools/gen1.py (auto); review=proofread-sonnet; flags=none
\begin{question}{46}
원점에 대해 오목한 생산가능곡선에 관한 설명으로 옳지 않은 것은?
\begin{choices}
  \item X축 상품생산이 늘어나면 기울기가 더 가팔라진다.
  \item 생산기술이 향상되면 생산가능곡선이 원점에서 더 멀어진다.
  \item 기회비용 체증의 법칙이 성립한다.
  \item 생산가능곡선 기울기의 절댓값이 한계변환율이다.
  \item 생산가능곡선 상의 점에서 파레토 개선이 가능하다.
\end{choices}
\end{question}
\answer{5}
